\subsection{실험 환경 및 결과}

Arduino Nano RP2040 Connect(125 MHz), arm-none-eabi-GCC 7.2(-O3)를 사용하였고, 시간은
1 MHz 하드웨어 타이머로 측정하였다. 표~\ref{tab:result}는 원소당 평균 처리 시간이다.

\begin{table}[t]
\centering
\caption{50,000원소 처리 성능 (원소당 평균)}
\label{tab:result}
\small
\begin{tabular}{lrrr}
\toprule
구현 & ns/op & 사이클 & 기준 대비 \\
\midrule
origin (분기) & 120.2 & 15.0 & 1.00$\times$ \\
shift\_pow2 (branchless C) & 144.3 & 18.0 & 0.83$\times$ \\
\textbf{sum\_dual\_asm} & \textbf{80.2} & \textbf{10.0} & \textbf{1.50$\times$} \\
\bottomrule
\end{tabular}
\end{table}

분기 제거 C 구현은 기준보다 오히려 1.2배 느렸고, 같은 목적의 수동 어셈블리는 1.5배
빨랐다. 원인을 확인하기 위해 컴파일러가 실제로 생성한 코드를 분석하였다.

\subsection{컴파일러 생성 코드와의 차이 분석}

GCC -O3가 shift\_pow2 루프에 생성한 실제 코드는 원소당 14명령어다(루프 제어 포함).

\begin{lstlisting}
movs r3,#0        @ zero offset
movs r4,r6        @ copy const 31
movs r7,r5        @ copy const 1
ldrsb r3,[r2,r3]
adds r2,r2,#1
adds r3,r3,#5     @ e = s+5
ands r4,r3        @ e & 31
lsls r7,r7,r4     @ 1 << (e&31)
movs r4,r7        @ copy result
asrs r3,r3,#31    @ sign mask
bics r4,r3        @ apply mask
adds r0,r0,r4
cmp r1,r2
bne .L3
\end{lstlisting}

수동 어셈블리(8명령어)와의 차이는 세 층위로 나뉜다.

첫째, \textbf{언어 규칙에 의한 알고리즘 제약}이다. 포화 특성을 쓰면 부호 처리가 0명령어인데,
C의 미정의 동작 규칙 때문에 컴파일러는 이 변환 자체가 금지되어 있다. 그래서 branchless C는
마스크 생성--적용(\texttt{asrs}, \texttt{bics})과 시프트 양 클램프(\texttt{ands})라는 우회
경로를 강제당한다. 이것은 최적화기의 능력 문제가 아니라 규칙 문제다.

둘째, \textbf{명령어 집합 구조에 의한 낭비}다. Thumb-1은 대부분 2피연산자 파괴형 명령어라
상수(1, 31)를 보존하려면 매 반복 레지스터 복사(\texttt{movs r4,r6} 등)가 필요하다. 위
코드에서 원소당 3개의 \texttt{movs}가 순수한 이동 비용이다.

셋째, \textbf{이론과 실측의 오차가 작다}는 점이다. 수동 어셈블리의 이론 사이클은 명령어
8개 + 로드 추가 1 + 루프 분담 1.25 = 10.25인데 실측은 10.0으로 오차 2\% 이내다. M0+의
결정적 실행 모델에서는 명령어 수 분석만으로 성능을 정확히 예측할 수 있으며, 반대로 말하면
컴파일러가 만든 14명령어와 사람이 만든 8명령어의 차이가 그대로 성능 차(18 대 10사이클)로
나타난다.

한편 동일한 어셈블리 아이디어라도 조각(fragment) 단위로 C 루프에 삽입하면 기준과 동률
(120.2 ns)에 그쳤다. 컴파일러는 인라인 어셈블리 내부를 볼 수 없어 조각마다 레지스터 4개를
통째로 예약하고, -O3 언롤 과정에서 8개뿐인 low 레지스터가 고갈되어 스택 스필이 발생하기
때문이다. 어셈블리 최적화는 조각이 아니라 루프 단위로 적용해야 효과가 있다.

\subsection{논의 및 한계}

이 결과에서 세 가지 교훈을 끌어낼 수 있다. 첫째, 분기 예측기가 없는 코어에서 분기는 이미
1--3사이클로 저렴하므로, 분기 하나를 없애기 위해 ALU 연산을 네 개 이상 추가하는 변환은
손해다. branchless가 유리하다는 상식은 마이크로아키텍처에 따라 역전된다. 둘째, 하드웨어가
제공하는 유용한 동작(시프트 포화)이 언어 표준에서는 미정의로 묶여 있을 때, 그 간극은 컴파일러
버전이 올라가도 좁혀지지 않는다. 규칙의 문제이지 구현의 문제가 아니기 때문이다. 셋째,
결정적 실행 코어에서는 종이 위의 사이클 계산이 실측과 2\% 이내로 맞으므로, 최적화 여지가
있는지 여부를 코드를 실행하기 전에 판정할 수 있다.

본 실험의 한계도 명확히 해 둔다. 수동 어셈블리 구현은 값이 상수 32라는 조건과 시프트
지시자가 $\pm 31$ 이내라는 조건을 전제한다. 값이 배열로 주어지면 상수 적재
(\texttt{movs})가 로드(\texttt{ldr})로 바뀌어 원소당 2--3사이클이 추가되지만 포화 트릭
자체는 그대로 유효하다. 또한 본 결과는 분기 예측기가 있는 상위 코어(Cortex-M7 등)로
일반화할 수 없으며, 오히려 그런 코어에서는 branchless가 다시 유리해질 수 있다. 즉 최적
구현은 코어마다 다시 실측되어야 하고, 이것이 다음 장에서 논의할 자동화가 필요한 이유다.
